% Part: applied-modal-logics
% Chapter: temporal-logic
% Section: possible-histories

\documentclass[../../../include/open-logic-section]{subfiles}

\begin{document}

\olfileid{aml}{tl}{poss}

\olsection{સંભવિત ઇતિહાસો}

અત્યાર સુધી વાપરેલાં સમયસંબંધી તર્કનાં સંબંધાત્મક નિદર્શનો ઘણાં
લવચીક છે, કારણ કે સુલભતા સંબંધ પર કોઈ નિયંત્રણ મૂકવું પડતું નથી.
એટલે સમયસંબંધી નિદર્શનોમાં ભૂતકાળ અને ભવિષ્ય બંને તરફ શાખાઓ પડી
શકે છે. પરંતુ આપણે શાખાઓનો વધુ ``મોડલ'' ખ્યાલ લઈ શકીએ, જેમાં
ઘટનાઓની શ્રેણીઓને સંભવિત ઇતિહાસો ગણીએ. આ માટે ભાષા બદલવી જરૂરી
નથી; જોકે આપણે આપણા ``સામાન્ય'' મોડલ સંક્રિયકો~$\Box$ અને~$\Diamond$
પણ ઉમેરી શકીએ, તથા સમય સાથે કર્તાઓના જ્ઞાનમાં થતા ફેરફારો દર્શાવવા
જ્ઞાનસંબંધી સુલભતા સંબંધો ઉમેરવાનું પણ વિચારી શકીએ.

\begin{defn}
  સમયસંબંધી ભાષાનું \emph{સંભવિત ઇતિહાસોનું નિદર્શ} એ ત્રિક
  $\mModel{M} = \tuple{T, C, V}$ છે, જ્યાં
  \begin{enumerate}
    \item $T$ સમયની સ્થિતિઓ તરીકે અર્થઘટિત થતો અખાલી ગણ છે.
    \item $C$ એ તંત્રના સંગણનાત્મક પથોનો, એટલે કે \emph{સંભવિત
      ઇતિહાસો}નો ગણ છે. બીજા શબ્દોમાં, $C$ એ સ્થિતિઓ
      $s_1$, $s_2$,~$s_3$, \dots ની શ્રેણીઓ~$\sigma$નો ગણ છે,
      જ્યાં દરેક $s_i \in T$ છે.
    \item $V$ એવું વિધેય છે જે દરેક !!{propositional
      variable}~$p$ને સમયબિંદુઓનો ગણ~$V(p)$ આપે છે.
  \end{enumerate}
  સરળતા માટે સામાન્ય રીતે એવું પણ માનીએ છીએ કે કોઈ ઇતિહાસ~$C$માં
  હોય, તો તેની બધી અંતિમ ઉપશ્રેણીઓ પણ તેમાં હોય. દાખલા તરીકે,
  $s_1$, $s_2$,~$s_3$ શ્રેણી~$C$માં હોય, તો $s_2$, $s_3$ અને~$s_3$
  પણ તેમાં હોય. વધુમાં, બે સ્થિતિઓ $s_i$ અને~$s_j$ શ્રેણી~$\sigma$માં
  આવે ત્યારે, $i < j$ હોય તો $s_i \prec_\sigma s_j$ કહીએ છીએ.
  $t \in V(p)$ હોય ત્યારે કહીએ કે $p$ એ~$t$એ \emph{સત્ય છે}.
\end{defn}

અહીં સંબંધિત એક ફેરફાર એ છે કે નિદર્શ~$\mModel M$માં
સમયબિંદુ~$t$એ !!a{formula}નું સત્ય ચકાસીએ ત્યારે તેને એવા
ઇતિહાસ~$\sigma$ને અનુલક્ષીને ચકાસીએ છીએ જેમાં~$t$ એક સ્થિતિ તરીકે
આવે છે. જોકે !!{propositional
variable}s અથવા સત્યવિધેયી
સંયોજકોનું અર્થવિજ્ઞાન બદલવાની જરૂર નથી. તેઓ~$\sigma$નો ઉલ્લેખ
કરતા નથી, એટલે તેમના નિયમો \olref[sem]{defn:tmodels} જેવા જ છે.
પરંતુ હવે આ ઇતિહાસોને અનુલક્ષીને ભવિષ્ય-સંક્રિયક~$\Ftemp$ને ફરી
વ્યાખ્યાયિત કરીએ છીએ અને સંક્રિયક~$\Diamond$ ઉમેરીએ છીએ.

\begin{defn}\ollabel{defn:phmodels}
  નિદર્શ~$\mModel M$માં \emph{!!a{formula}~$!A$નું~$t, \sigma$એ સત્ય હોવું},
  સંજ્ઞામાં $\mSat{M}{!A}[t, \sigma]$:
  \begin{enumerate}
  \item\ollabel{defn:sub:phmodels-f} \indcase{!A}{\Ftemp
    !B}{$\mSat{M}{\indfrm}[t, \sigma]$ ત્યારે અને માત્ર ત્યારે જ જ્યારે
    $t \prec_\sigma t'$ ધરાવતા કોઈ $t' \in T$ માટે
    $\mSat{M}{!B}[t', \sigma]$.}
  \item\ollabel{defn:sub:phmodels-diamond} \indcase{!A}{\Diamond
    !B}{$\mSat{M}{\indfrm}[t, \sigma]$ ત્યારે અને માત્ર ત્યારે જ જ્યારે
    $t$ આવતું હોય એવા કોઈ $\sigma' \in C$ માટે
    $\mSat{M}{!B}[t, \sigma']$.}
  \end{enumerate} 
\end{defn}

બીજા સમયસંબંધી અને મોડલ સંક્રિયકો પણ એ જ રીતે વ્યાખ્યાયિત કરી
શકાય. પરંતુ હવે કાળ અને શક્યતાને જોડતા દાવાઓ વ્યક્ત કરી શકીએ છીએ.
દાખલા તરીકે, ``$p$~બનશે નહીં, પણ તે બની શક્યું હોત'' એ વાત
$\lnot \Ftemp p \land \Diamond \Ftemp p$ સૂત્રથી દર્શાવી શકીએ.
આ સૂત્ર એવા બિંદુ અને ઇતિહાસમાં સત્ય હશે જ્યાં પછીની કોઈ પણ સ્થિતિએ
$p$ સત્ય થતું નથી, પરંતુ કોઈ વૈકલ્પિક ઇતિહાસમાં $p$ આગળ સત્ય થશે.

\end{document}
